% Figure: a friction wedge driven under a load, with the FBD of the wedge
% showing the two contact reactions and the driving force at the verge of
% slipping. Supports the limiting-equilibrium worked example.
\begin{figure}[htbp]
\centering
\begin{tikzpicture}[
  force/.style={draw=engblue, -{Stealth}, very thick},
  nrm/.style={draw=engred, -{Stealth}, very thick},
  fric/.style={draw=engamber, -{Stealth}, very thick},
  scale=1.0,
]
% ===== Left: the scene =====
\begin{scope}[xshift=0cm]
  % ground
  \draw[draw=engblack, very thick] (-0.4,0) -- (5.0,0);
  \foreach \x in {-0.2,0.2,...,4.8} {\draw[draw=enggray, thin] (\x,0) -- (\x-0.2,-0.25);}
  % the block being lifted (rests on top of the wedge incline)
  \draw[draw=engblack, fill=engblue!10, thick] (2.4,0.9) rectangle (4.6,2.1);
  \node[font=\scriptsize] at (3.5,1.5) {load $W$};
  % the wedge: thin triangle, apex angle alpha, driven to the right under
  % the block. base on the ground, top face under the block.
  \draw[draw=engblack, fill=engamber!15, thick]
    (0.2,0) -- (4.6,0) -- (4.6,0.9) -- cycle;
  \node[font=\scriptsize] at (3.0,0.32) {wedge};
  % apex angle alpha at the thin (left) end
  \draw[draw=enggray, thin] (0.9,0) arc (0:11:0.7);
  \node[font=\scriptsize, anchor=west] at (0.95,0.08) {$\alpha$};
  % driving force P on the wedge, horizontal to the right
  \draw[force] (0.2,0.18) -- ++(-1.2,0)
    node[anchor=east, font=\small, engblue] {$P$};
  \node[font=\small, anchor=north] at (2.2,-0.45) {(a) the scene};
\end{scope}
% ===== Right: FBD of the wedge alone =====
\begin{scope}[xshift=7.0cm]
  % wedge outline repeated
  \draw[draw=engblack, fill=engamber!15, thick]
    (0.2,0) -- (4.6,0) -- (4.6,0.9) -- cycle;
  % driving force
  \draw[force] (0.2,0.18) -- ++(-1.3,0)
    node[anchor=east, font=\scriptsize, engblue] {$P$};
  % bottom contact (with ground): normal up, friction left (opposing the
  % wedge's rightward motion)
  \draw[nrm] (2.6,0) -- ++(0,1.3) node[anchor=south, font=\scriptsize, engred] {$N_1$};
  \draw[fric] (2.6,0) -- ++(-1.3,0) node[anchor=south, font=\scriptsize, engamber] {$F_1$};
  % top contact (with the load): normal perpendicular to the inclined top
  % face, pointing down-and-left into the wedge; friction along that face
  % opposing the wedge sliding under the block (down-the-incline on wedge)
  \draw[nrm] (3.6,0.74) -- ++(-0.5,-1.25)
    node[anchor=north, font=\scriptsize, engred] {$N_2$};
  \draw[fric] (3.6,0.74) -- ++(-1.25,-0.25)
    node[anchor=east, font=\scriptsize, engamber] {$F_2$};
  \node[font=\small, anchor=north] at (2.2,-0.45) {(b) FBD of the wedge};
\end{scope}
\end{tikzpicture}
\caption[A friction wedge at limiting equilibrium]{A wedge of apex angle
$\alpha$ driven horizontally under a load $W$. (a) The scene: the driving
force $P$ pushes the wedge to the right, lifting the load along the
wedge's top face. (b) The free-body diagram of the wedge has two
contacts, each contributing a normal force and a friction force. At the
ground the normal $N_1$ is vertical and the friction $F_1$ opposes the
rightward motion. At the top face the normal $N_2$ is perpendicular to the
incline and the friction $F_2$ acts along it, both opposing the wedge
sliding under the block. At the verge of motion each friction takes its
limiting value $\mu N$, and the three equilibrium equations close on the
driving force $P$ needed to advance the wedge.}
\label{fig:vol03ch01:wedge}
\end{figure}
